
This study tested whether adversarial fragility and calibration error are positively coupled in large language models after controlling for general capability. Across 30 LLMs spanning 9 families, the Residual Instability (RI) construct — the OLS residual of AdvGLUE accuracy drop after regressing out a PCA-derived capability index and mean confidence — significantly anticorrelates with Expected Calibration Error on arc_challenge (Spearman partial ρ = −0.535, p = 0.0034, 95\% CI = [−0.782, −0.101]). This direction is opposite to the pre-registered prediction of ρ ≥ +0.4.

Three contributions are reported: (1) the RI construct, validated as non-degenerate (R² = 0.529, SD = 0.121, VIF = 1.000); (2) the empirical anticorrelation between RI and ECE, which refutes the coupled failure cascade prediction for this dimension pair; and (3) a calibration–robustness trade-off hypothesis, grounded in the RLHF literature, as the most plausible mechanistic explanation for the observed anticorrelation.

Several limitations constrain the strength of these conclusions. Most critically, 73\% of AdvGLUE scores are OLS-estimated rather than directly measured, and ECE is derived from a single benchmark. The study scope is limited to the RI–ECE dimension; sub-hypotheses for hallucination, safety, and output variance were not executed. Future work should include direct AdvGLUE measurement on all 30 models, ECE replication across multiple benchmarks, and training-regime stratification to test the proposed trade-off mechanism. Whether the inverted relationship extends to hallucination and safety dimensions requires executing redesigned sub-hypotheses with revised directional predictions.

The result does not support treating adversarial fragility and calibration as positively coupled failure modes in the context studied. Multi-dimension trustworthiness evaluations may benefit from empirically testing cross-dimension relationships rather than assuming a particular coupling structure.

